{\namefont\bfseries\color{accentgray}Баранов Константин}\hfill {\rolefont\bfseries\color{accentgray}C++ Software Engineer}\par
\vspace{1.6pt}
{\color{mutedgray}3+ года профессионального опыта\sep Санкт-Петербург, Россия\sep Офис / гибрид / удалённо\sep Готов к релокации}\par
\vspace{0.5pt}
{\contactfont \href{tel:+79111747978}{+7 (911) 174-79-78}\sep \href{mailto:job@seig.ru}{job@seig.ru}\sep \href{https://t.me/seigtm}{Telegram: @seigtm}\sep \href{https://www.linkedin.com/in/seigtm}{linkedin.com/in/seigtm}}\par
{\contactfont \href{https://github.com/seigtm}{github.com/seigtm}\sep \href{https://seigtm.github.io}{seigtm.github.io}\sep Русский: родной\sep Английский: C2 Proficient, EF SET 71/100 (\href{https://cert.efset.org/du7yCw}{cert.efset.org/du7yCw})}\par

\section*{Технические навыки}
\skill{Языки}{C++20, C, Python, Bash, PowerShell.}
\skill{C++ / сборка / инструменты}{STL, CMake, Android NDK, JNI, Conan, Google Test, CTest, Git, GitLab CI, clangd, \mbox{clang-format}, Clang, GCC, MSVC.}
\skill{Платформы и системная разработка}{Linux, Windows, Android, POSIX API, IPC, клиент-серверная архитектура, Protobuf.}
\skill{Android / reverse engineering}{Android internals, Binder service calls, boot control, init.rc, ramdisk, fstab, контексты SELinux, Magisk, smali/baksmali, JADX, IDA.}
\skill{Графика / desktop}{OpenGL, SDL, Dear ImGui, Assimp, GLSL, PySide6.}

\section*{Профессиональный опыт}

\role{C++ Software Engineer}{декабрь 2022 -- настоящее время}{ООО «Специальный Технологический Центр»}{Санкт-Петербург, Россия}
\begin{itemize}
  \item Веду полный цикл разработки нативного C/C++ ПО для Android, Linux и Windows: от архитектурного планирования и реализации внутренних инженерных проектов до демонстрации и защиты программных комплексов перед крупными заказчиками.
  \item Разработал нативный C++20-сервис для Android с IPC-интерфейсом и клиент-серверным API для управления системными функциями устройства.
  \item Проектировал обмен данными между Android/Linux-компонентами, используя TCP, локальные IPC-каналы и Protobuf для взаимодействия клиентских и серверных частей.
  \item Настраивал Android NDK/CMake-сборку и GitLab CI для сборки под несколько целевых платформ и публикации артефактов в корпоративное облачное хранилище.
  \item Работал с Android internals: boot control, Binder/telephony service calls и privileged execution contexts; модифицировал системные образы Android, ramdisk, init/fstab и контексты SELinux.
  \item Анализировал и модифицировал Android/Linux/Windows-приложения, библиотеки и системные компоненты с использованием smali, baksmali, JADX и IDA.
  \item Проводил технические исследования проприетарных форматов сообщений модемов Qualcomm/MediaTek, сетевых протоколов, Bluetooth и спецификаций 3GPP.
  \item Участвую в code review и инициирую обсуждение архитектурных решений.
\end{itemize}

\role{Junior C++ Software Developer}{сентябрь 2021 -- январь 2022}{ООО «VAS Experts»}{Санкт-Петербург, Россия}
\begin{itemize}
  \item Разрабатывал C++-компоненты распределённой системы хранения и передачи файлов.
  \item Реализовывал обработку сетевых команд, Protobuf-сериализацию и преобразование данных между внутренними протоколами системы.
  \item Проектировал конечные автоматы для управления состояниями сессий, файлов и хранилищ; покрыл модульными тестами на Google Test ключевые модули и критичные взаимодействия.
  \item Настраивал CMake-сборку и эксплуатационную инфраструктуру Linux-сервисов: модульная структура проекта, CTest, горячая перезагрузка конфигурации и ротация логов.
\end{itemize}

\role{Intern Backend Developer}{июнь 2021}{ГУП «Санкт-Петербургский информационно-аналитический центр»}{Санкт-Петербург, Россия}
\begin{itemize}
  \item Разрабатывал вспомогательные инструменты и backend-компоненты для внутренних задач команды с использованием окружений Docker и PostgreSQL.
\end{itemize}

\role{Intern C++ Software Developer}{май 2021}{АО «Невское проектно-конструкторское бюро»}{Санкт-Петербург, Россия}
\begin{itemize}
  \item Разработал desktop-приложение на C++/Qt от проектирования до внедрения.
\end{itemize}

\section*{Проекты}

\project{MEOV -- Minimalistic Easy Object Viewer}{\href{https://github.com/seigtm/meov}{github.com/seigtm/meov}}
\begin{itemize}
  \item Создал desktop-приложение для просмотра 3D-моделей на C++ с использованием OpenGL, SDL, Dear ImGui, Assimp и CMake.
  \item Спроектировал компонентную архитектуру сцены с поддержкой моделей, материалов, текстур, шейдеров, камеры, освещения, skybox и dockable UI.
  \item Настроил CMake-сборку, статическую линковку зависимостей, Git submodules, runtime assets и настройки sanitizer/debug.
  \item Реализовал ресурсный слой с кешированием моделей, текстур и шейдерных программ; интегрировал загрузку meshes, materials и texture maps через Assimp.
  \item Защитил проект как бакалаврскую выпускную квалификационную работу с отличием.
\end{itemize}

\section*{Образование}

\textbf{Санкт-Петербургский политехнический университет Петра Великого}\hfill 2022 -- 2025\par
Бакалавриат, 09.03.04 Программная инженерия; Институт компьютерных наук и кибербезопасности\par
\tightvspace
\textbf{Санкт-Петербургский политехнический колледж городского хозяйства}\hfill 2018 -- 2022\par
Программирование в компьютерных системах; техник-программист и оператор ЭВМ, дипломы с отличием
